\documentclass[10pt,a4paper]{amsart}
\usepackage[margin=19mm]{geometry}
\usepackage{amsmath,amssymb,mathtools}
\usepackage[scr=rsfs,cal=euler]{mathalfa}
\usepackage{xfrac}
\usepackage[unicode,hidelinks,bookmarks=false]{hyperref}
\newcommand{\ie}{i.e.\ }
\newcommand{\cf}{cf.\ }
\newcommand{\resp}{respectively\ }
\newcommand{\Subsection}{\subsection}
\providecommand{\nobreakdash}{\nobreak\hbox{-}}
\setlength{\emergencystretch}{2em}
\multlinegap0pt
\usepackage{enumerate}
\usepackage{amssymb}
\usepackage{cancel}
\newtheorem{lemma}{Lemma}
\renewcommand{\div}{\,\textup{div}}
\newcommand{\lot}{\textup{l.o.t.}}
\newcommand{\tr}{\textup{tr}}
\title{A Spinorial Perelman's Functional: critical points and gradient flow}
\author{Tsz-Kiu Aaron Chow, Frederick Tsz-Ho Fong}
\date{Source: arXiv:2601.01863v1 (CC BY 4.0)}
\begin{document}
\maketitle

\noindent\emph{Source and licence.} A Spinorial Perelman's Functional: critical points and gradient flow. Version arXiv:2601.01863v1 (CC BY 4.0), distributed by arXiv under the Creative Commons Attribution 4.0 International licence, http://creativecommons.org/licenses/by/4.0/, as recorded in the arXiv licence metadata for this exact version and shown on the abstract page https://arxiv.org/abs/2601.01863v1. Full licence notice: \texttt{licenses/arxiv-2601.01863-CC-BY-4.0.txt}.\par\smallskip

\noindent\textit{Editorial scope.} Counted unit: the article's lemma lem;symbol\_2 (source0.tex lines 861--874). The article's complete original proof of it is delivered with it (source0.tex lines 876--947, one \texttt{proof} environment of 72 lines). No mathematics is written, repaired or re-derived: the statement and the proof are the article's own lines, in the article's own notation, and no retained line is substituted or re-flowed.  No further source-local context is needed: every cross-reference of every retained line resolves inside the two delivered spans.  Nothing was omitted from any retained span, so the package carries no omission record.  No retained line contains a citation, so the wrapper carries no bibliography and no bibliography entry.  No retained line contains a figure environment, an image-inclusion command or drawing code, so no image is delivered and the package declares no asset dependency.  Editorial formatting only, in addition: an AMS \texttt{amsart} preamble that restates the article's own declarations and macros used by the retained lines, namely the environments \texttt{lemma} on separate counters, together with the standard packages those lines need.  The retained environments print in this document as Lemma 1 in order of appearance, whereas the article prints the same items with the numbering of its own full document.  The complete native source of the article as distributed in the arXiv e-print for this exact version is bundled unchanged as \texttt{sources/192114/source0.tex}.
\begin{lemma}\label{lem;symbol_2}
	Let $(\dot{g}, \dot{\psi}) = (\eta, s)$. We have
	\begin{align*}
		\sigma_{\xi}\left(  D_{(g,\psi)} \mathcal{L}_{U}^{\text{spin}}\psi \right)(\eta,s)
		&  = -\frac{1}{16}\sum_{i,a,b}\sum_{j\neq k}\, \xi_a\xi_j\, \eta(e_i, e_k)\, \textup{Re}\langle \psi,\, e_b\cdot e_i\cdot e_j\cdot e_k\cdot\psi\rangle \, e_a\cdot e_b\cdot \psi\\[1ex]
		&\quad  + \frac{1}{16}\sum_{i,a,b}\sum_{j\neq k}\, \xi_b\xi_j\, \eta (e_i, e_k)\, \textup{Re}\langle \psi,\, e_a\cdot e_i\cdot e_j\cdot e_k\cdot\psi\rangle \, e_a\cdot e_b\cdot\psi\\[1ex]
		&\quad -  \frac{1}{4}\sum_{a, b}\xi_b\, \textup{Re}\langle\psi,\, e_a\cdot \xi\cdot s\rangle \, e_a\cdot e_b\cdot \psi +  \frac{1}{4}\sum_{a, b} \xi_a\, \textup{Re}\langle\psi,\, e_b\cdot \xi\cdot s\rangle  \, e_a\cdot e_b\cdot \psi,\\[1ex]
		\sigma_{\xi}\left( D_{(g,\psi)} \mathcal{L}_{W(g,g_0) }^{\text{spin}}\psi \right)(\eta,s) &= -\frac{1}{4} \sum_k\sum_{i\neq j}(\xi_i\xi_k\,\eta(e_j, e_k) - \xi_j\xi_k\eta(e_i, e_k))e_i\cdot e_j\cdot \psi,\\[1ex]
	\sigma_{\xi}\left( D_{(g,\psi)}\Delta_g\psi \right)(\eta,s) &= -\frac{1}{4}\sum_i\sum_{j\neq k}\xi_i\xi_j\, \eta(e_i,e_k)e_j\cdot e_k\cdot\psi - |\xi|^2 s,\\[1ex]
	\sigma_{\xi}\left( D_{(g,\psi)}\div_g\, U \right)(\eta,s) &=  \frac{1}{4}\sum_i\sum_{j\neq k}\xi_i\xi_j\, \eta(e_i,e_k)\, \text{Re}\langle\psi,\,e_j\cdot e_k\cdot\psi\rangle + |\xi|^2\,\text{Re}\langle\psi, s\rangle\\
	&\quad  + \frac{1}{4}|\xi|^2\tr(\eta)\,|\psi|^2 -  \frac{1}{4}\eta(\xi,\xi)\,|\psi|^2.
	\end{align*}

\end{lemma}

\begin{proof}
	Since we wish to compute the principal symbol at the end, we only need to worry about the linearization at top order.	We have
	\begin{align}
		\left(D_{(g,\psi)}\mathcal{L}_{U}^{\text{spin}}\psi\right)(\eta,s) & = \mathcal{L}_{(D_g U)(\eta)}^{\text{spin}}\psi  + \mathcal{L}_{(D_\psi U)(s)}^{\text{spin}}\psi + \lot
	\end{align}
	
	Note that $U^a = \text{Re}\langle\psi, e_a\cdot D\psi\rangle $. 
	Using [AWW, Lemma 4.12], we calculate, up to highest order, 
	\begin{align}
		(D_g U^a)(\eta) &= \frac{1}{4}\sum_{i}\sum_{j\neq k}\, (\nabla_{e_j}\eta)(e_i, e_k)\, \textup{Re}\langle \psi,\, e_a\cdot e_i\cdot e_j\cdot e_k\cdot\psi\rangle  + \lot
	\end{align} 
	This gives
	\begin{equation}
		\begin{aligned}
			\mathcal{L}_{(D_g U)(\eta)}^{\text{spin}}\psi &=  -\frac{1}{4}\sum_{a\neq b}(\nabla_a (D_g U^b(\eta)) - \nabla_b (D_g U^a(\eta)))\, e_a\cdot e_b\cdot \psi + \lot\\
		&= -\frac{1}{16}\sum_{i,a,b}\sum_{j\neq k}\, (\nabla_{e_a}\nabla_{e_j}\eta)(e_i, e_k)\, \textup{Re}\langle \psi,\, e_b\cdot e_i\cdot e_j\cdot e_k\cdot\psi\rangle\, e_a\cdot e_b\cdot \psi\\
		&\quad + \frac{1}{16}\sum_{i,a,b}\sum_{j\neq k}\, (\nabla_{e_b}\nabla_{e_j}\eta)(e_i, e_k)\, \textup{Re}\langle \psi,\, e_a\cdot e_i\cdot e_j\cdot e_k\cdot\psi\rangle\, e_a\cdot e_b\cdot \psi + \lot
		\end{aligned}
	\end{equation}
For the another other top order terms\, we also calculate, up to highest order,
	\begin{align}
		&(D_\psi U^a)(s) = \textup{Re}\langle\psi,\, e_a\cdot Ds\rangle + \lot
	\end{align}
	This gives
	
	\begin{equation}
		\begin{aligned}
			\mathcal{L}_{(D_{\psi} U)(s)}^{\text{spin}}\psi &=  -\frac{1}{4}\sum_{a\neq b}(\nabla_a (D_{\psi} U^b(s)) - \nabla_b (D_{\psi}  U^a(s)))\, e_a\cdot e_b\cdot \psi + \lot\\
			&= -\frac{1}{4}\sum_{a\neq b}\textup{Re}\langle\psi,\, e_a\cdot \nabla_{e_b}Ds\rangle \, e_a\cdot e_b\cdot \psi  + \frac{1}{4}\sum_{a\neq b}\textup{Re}\langle\psi,\, e_b\cdot \nabla_{e_a}Ds\rangle \, e_a\cdot e_b\cdot \psi + \lot\\
		\end{aligned}
	\end{equation}
Consequently,

	\begin{align}
		&\sigma_{\xi}\left(D_{(g,\psi)}\mathcal{L}_{U}^{\text{spin}}\psi\right)(\eta,s)\\
		&= -\frac{1}{16}\sum_{i,a,b}\sum_{j\neq k}\, \xi_a\xi_j\, \eta(e_i, e_k)\, \textup{Re}\langle \psi,\, e_b\cdot e_i\cdot e_j\cdot e_k\cdot\psi\rangle \, e_a\cdot e_b\cdot \psi\\[1ex]
		&\quad  + \frac{1}{16}\sum_{i,a,b}\sum_{j\neq k}\, \xi_b\xi_j\, \eta (e_i, e_k)\, \textup{Re}\langle \psi,\, e_a\cdot e_i\cdot e_j\cdot e_k\cdot\psi\rangle \, e_a\cdot e_b\cdot\psi\\[1ex]
		&\quad -  \frac{1}{4}\sum_{a, b}\xi_b\, \textup{Re}\langle\psi,\, e_a\cdot \xi\cdot s\rangle \, e_a\cdot e_b\cdot \psi +  \frac{1}{4}\sum_{a, b} \xi_a\, \textup{Re}\langle\psi,\, e_b\cdot \xi\cdot s\rangle  \, e_a\cdot e_b\cdot \psi.
	\end{align}


Secondly,
\begin{align*}
	D_{(g,\psi)}(\mathcal{L}_{W(g,g_0)}^{\text{spin}}\psi)(\eta,s) =  -\frac{1}{4}\sum_{i\neq j}(\nabla_i (D_g W^j(\eta)) - \nabla_j (D_g W^i(\eta)))\, e_i\cdot e_j\cdot \psi + \lot,
\end{align*}
This gives
\begin{align*}
	\sigma_{\xi}(D_{(g,\psi)}(\mathcal{L}_{W(g,g_0) }^{\text{spin}}\psi))(\eta,s) &= -\frac{1}{4}\sum_k\sum_{i\neq j}(\xi_i\xi_k\,\eta(e_j, e_k) - \xi_j\xi_k\eta(e_i, e_k))e_i\cdot e_j\cdot \psi.
\end{align*}


Next, using [AWW, Lemma 4.12] again, we have
\begin{align*}
	D_{(g,s)}(\Delta_g\psi)(\eta,s) &= \frac{1}{4}\sum_i\sum_{j\neq k}(\nabla_{e_i}\nabla_{e_j}\eta)(e_i,e_k)e_j\cdot e_k\cdot\psi + \Delta_g s + \lot
\end{align*}
This gives
\begin{align*}
	\sigma_{\xi}(D_{(g,s)}(\Delta_g\psi))(\eta,s) &= -\frac{1}{4}\sum_i\sum_{j\neq k}\xi_i\xi_j\, \eta(e_i,e_k)e_j\cdot e_k\cdot\psi - |\xi|^2 s.
\end{align*}



Lastly, note that 
\begin{align*}
	\div_g\, U = \text{Re}\langle\psi, D^2\psi\rangle= -\text{Re}\langle\psi, \Delta\psi\rangle + \frac{1}{4}R \,|\psi|^2
\end{align*}
Therefore, using the symbols for $\Delta_g\psi$ and $R_g$, we obtain
\begin{align*}
	\sigma_{\xi}(D_{(g,s)}(\div_g\, U))(\eta,s) &=  \frac{1}{4}\sum_i\sum_{j\neq k}\xi_i\xi_j\, \eta(e_i,e_k)\, \text{Re}\langle\psi,\,e_j\cdot e_k\cdot\psi\rangle + |\xi|^2\,\text{Re}\langle\psi, s\rangle\\
	&\quad + \frac{1}{4}|\xi|^2\tr(\eta)\,|\psi|^2 -  \frac{1}{4}\eta(\xi,\xi)\,|\psi|^2.
\end{align*}
\end{proof}

\end{document}
